% SUBSTITUTE for the "Jacobs Landscape Poster" template (originally
% bardsoftware / Vel, CC-BY-NC-SA — non-commercial, so it cannot be
% redistributed here).
%
% This starter uses the upstream `tikzposter` document class (CTAN,
% LPPL-1.2) in LANDSCAPE orientation on A0 paper, matching the
% original's purpose: a large-format academic conference poster. See
% README.md for the full disclosure.
%
% tikzposter is not bundled: it ships in every modern TeX distribution.
\documentclass[a0paper, landscape]{tikzposter}
\usepackage{booktabs}

\usetheme{Wave}
\usecolorstyle[colorPalette=BlueGrayOrange]{Britain}

\title{Predicting Urban Heat Islands from Satellite Land-Cover Data}
\author{J. Alvarez, M. Kowalski, and R. Chen}
\institute{Department of Environmental Engineering, Coastal State University}

\begin{document}
\maketitle

\begin{columns}

\column{0.33}

\block{Motivation}{
Cities can be several degrees warmer than surrounding rural areas, an
effect known as the urban heat island (UHI). Fine-grained UHI maps are
expensive to produce from ground sensors alone, but land-cover
classification from satellite imagery is cheap and updated regularly.
We ask whether land-cover fractions predict surface temperature well
enough to substitute for a sparse sensor network.
}

\block{Data}{
\begin{itemize}
\item Landsat 8 thermal band imagery, 30 m resolution, 42 cities.
\item Land-cover fractions (impervious surface, tree canopy, water)
  from the National Land Cover Database.
\item Ground-truth surface temperature from 1{,}100 weather stations.
\end{itemize}
}

\column{0.34}

\block{Method}{
We fit a gradient-boosted regression tree mapping land-cover fractions
and time of day to surface temperature, validated with 10-fold
cross-validation stratified by city.

\begin{itemize}
\item Features: impervious \%, canopy \%, water \%, elevation,
  distance to coastline.
\item Target: land surface temperature (LST), $^\circ$C.
\item Baseline: linear regression on impervious surface fraction alone.
\end{itemize}
}

\block{Results}{
\begin{center}
\begin{tabular}{lcc}
\toprule
Model & RMSE ($^\circ$C) & $R^2$ \\
\midrule
Linear (impervious only) & 3.8 & 0.41 \\
Gradient-boosted trees & 1.6 & 0.87 \\
\bottomrule
\end{tabular}
\end{center}
The full feature set roughly halves prediction error relative to the
impervious-surface-only baseline, and canopy fraction is the single
most important feature after impervious surface.
}

\column{0.33}

\block{Discussion}{
Land-cover fractions alone explain most of the variance in surface
temperature across cities of very different sizes and climates,
suggesting the model transfers reasonably well without per-city
recalibration. Coastal cities show the largest residual errors,
plausibly from unmodelled sea-breeze cooling.
}

\block{Conclusion}{
Satellite-derived land cover is a practical, low-cost substitute for
dense ground sensor networks when estimating city-scale heat island
intensity, particularly for identifying priority neighbourhoods for
tree-planting interventions.
}

\block{Contact}{
j.alvarez@coastalstate.edu \hfill Poster session 4B, Thursday 14:00
}

\end{columns}

\end{document}
